\section{Introduction}\label{sec:introduction}

The ionization problem asks how the outer electrons of an atom behave
as its nuclear charge increases.  The generalized ionization conjecture
predicts universal asymptotics for both ionization energies and radii
defined by a prescribed exterior electron mass.  We resolve its radius
part positively for the full correlated two-spin Schr\"odinger model,
using the conditional screening and propagation estimates of
\cite{OuterRadius}.

\subsection{Statement of the result}

In atomic units, for each integer $Z\ge1$ let
\begin{equation}\label{intro:hamiltonian}
 H_Z=\sum_{i=1}^{Z}\left(-\frac12\Delta_{x_i}-\frac{Z}{|x_i|}\right)
       +\sum_{1\le i<j\le Z}\frac1{|x_i-x_j|}
 \quad\hbox{on }\bigwedge^ZL^2(\R^3;\C^2).
\end{equation}
The operator is associated with its usual closed, bounded-below quadratic
form.  Let $\Psi_Z$ be any normalized ground state and let its spin-summed
one-electron density be
\begin{equation}\label{intro:density}
 \rho_{\Psi_Z}(x)
 =Z\sum_{\sigma_1,\ldots,\sigma_Z\in\{1,2\}}
   \int_{\R^{3(Z-1)}}
   |\Psi_Z(x,\sigma_1;x_2,\sigma_2;\ldots;x_Z,\sigma_Z)|^2
   \dd x_2\cdots\dd x_Z.
\end{equation}
For $Z=1$ the integral is absent.  Thus $\int\rho_{\Psi_Z}=Z$.
For integers $1\le m<Z$, define
\begin{equation}\label{intro:radius}
 R_m(\Psi_Z)=\inf\left\{r\ge0:
               \int_{|x|>r}\rho_{\Psi_Z}(x)\dd x\le m\right\}.
\end{equation}
This is a radius in the neutral atom: $m$ specifies the expected electron
mass outside the sphere.  The state itself still has $Z$ electrons.

\begin{theorem}[Generalized ionization conjecture: radii]\label{intro:main}
Set $b_{\mathrm{TF}}=(81\pi^2/2)^{1/3}$.  For every choice of normalized
neutral ground states $(\Psi_Z)_{Z\ge1}$ of \eqref{intro:hamiltonian},
\begin{align}
 \lim_{m\to\infty}m^{1/3}
       \left[\limsup_{\substack{Z\to\infty\\ Z>m}}R_m(\Psi_Z)\right]
       &=b_{\mathrm{TF}},\label{intro:limsup}\\
 \lim_{m\to\infty}m^{1/3}
       \left[\liminf_{\substack{Z\to\infty\\ Z>m}}R_m(\Psi_Z)\right]
       &=b_{\mathrm{TF}}.\label{intro:liminf}
\end{align}
All limits are over integers.  The inner limits hold $m$ fixed; no
convergence in $Z$ at a fixed $m$ is assumed.
\end{theorem}

The statement is uniform over all choices in each ground eigenspace.
Neither the state nor its density is assumed to be spherically symmetric.
Existence of these ground states, in the same form and spin convention,
follows from the strict binding statement in
\cite[\FoundationBinding]{OuterRadius}.

\subsection{History and scope}

The Thomas--Fermi approximation became a rigorous description of the
leading energy and scaled bulk density through the work of Lieb and
Simon \cite[Theorems~III.1 and~III.3]{LiebSimon1977}.
A neutral atom has bulk length scale
$Z^{-1/3}$, while the radius in \eqref{intro:radius}, at fixed $m$,
concerns only the fraction $m/Z$ of its mass.  Convergence of the bulk
density therefore leaves this outer-electron question unresolved.
For the full quantum model, Seco, Sigal, and Solovej obtained a
quantitative lower bound on the radius containing all but one expected
electron \cite[Theorem~3]{SecoSigalSolovej1990}.  This concerns the same
expected-tail notion of radius, but its lower bound still decays with $Z$
and does not determine a universal outer-radius law.

Solovej's unrestricted Hartree--Fock theorem established the universal
outer-radius asymptotic with the same constant as in
\cref{intro:main} \cite[Definition~1.2 and Theorem~1.5]{Solovej2003}.
Its proof controls screened potentials on successively larger exterior
regions \cite[Sections~10--13]{Solovej2003}, uniformly in the chosen
Hartree--Fock minimizer.  The generalized ionization conjecture for the
full Schr\"odinger model explicitly retains both upper and lower
large-$Z$ limits \cite[Section~3, Equation~(2)]{Solovej2016}.
The order of limits in \cref{intro:main} is precisely this formulation.
Sharp radius estimates were subsequently obtained for minimizers in
Thomas--Fermi--Dirac--von Weizs\"acker theory
\cite[Theorem~1.3]{FrankNamVanDenBosch2018} and for power density-matrix
functionals \cite[Theorem~3]{Kehle2017}.  These are results for distinct
approximate variational models, not for arbitrary correlated
Schr\"odinger ground states.

The limiting profile has its own earlier analytic theory.  Lieb and
Simon proved directional Sommerfeld asymptotics for neutral
Thomas--Fermi systems \cite[Theorem~IV.10]{LiebSimon1977}.
Solovej's later comparison estimate explicitly allows nonradial
solutions \cite[Lemma~4.4 and Remark~4.5]{Solovej2003}; the homogeneous
rigidity statement used here is a consequence of that estimate.
We give a short dilation proof in \cref{prof:rigidity}.  The additional
task is to obtain such a profile from correlated quantum ground states
and then recover their expected exterior mass.

For the full model, the foundation companion already bounds above and
below, uniformly in $Z$ and the neutral ground state, the radius leaving
one half of an electron in expected exterior mass
\cite[Theorem~1.2]{OuterRadius}.  The present theorem identifies the
sharp coefficient as the prescribed exterior mass tends to infinity.
The same companion supplies the conditional density comparisons and
positive intermediate subsolutions used for this passage.
Section~\ref{sec:inputs} defines their common observations
and smoothing conventions; Appendix~\ref{sec:interface} records the
provider results and the order of parameter choices.
The intermediate-scale compactness strategy also appears in the
energy companion \cite[\EnergyProfile]{EnergyPart}.
That is a shared method: the proof here uses only the foundation
estimates and establishes its own neutral-profile and expected-tail
arguments.

\subsection{Proof and reusable ingredients}

The limiting neutral Thomas--Fermi model explains both the rescaling and
the radius coefficient.  Its density $\varrho$ and screened field $F$
satisfy, away from the nucleus,
\begin{equation}\label{intro:constants}
 \varrho=k(F_+)^{3/2},\qquad \Delta F=4\pi\varrho,
 \qquad k=\frac{2^{3/2}}{3\pi^2},\qquad t_+=\max(t,0).
\end{equation}
Here $k$ corresponds to two spin states and kinetic operator $-\Delta/2$
\cite[\FoundationConstants]{OuterRadius}.  Since
$\Delta|x|^{-4}=12|x|^{-6}$, the positive radial homogeneous solution is
$F_*(x)=A_*|x|^{-4}$, with
\begin{equation}\label{intro:sommerfeld}
 A_*:=\left(\frac{12}{4\pi k}\right)^2=\frac{81\pi^2}{8},
 \qquad \rho_*(x):=kA_*^{3/2}|x|^{-6}\quad(x\ne0).
\end{equation}
Its exterior mass is
\[
 \int_{|x|>r}\rho_*(x)\dd x=4A_*r^{-3}
       =b_{\mathrm{TF}}^3r^{-3},
\]
so prescribing mass $m$ gives radius $b_{\mathrm{TF}}m^{-1/3}$.
These are equations for the limiting model, not identities for the
quantum density $\rho_{\Psi_Z}$.

To reach this model from a quantum state, we add small independent errors
to the electron positions and condition on the resulting unordered
observations.  Conditionally averaging smooth unit-mass packets around
the electron positions gives a density; subtracting its Coulomb potential
from the nuclear potential gives the screened field.  At an intermediate
radius $s$, we rescale this field by $s^4$ and the density by $s^6$.
The comparison estimate gives a pointwise Thomas--Fermi relation on
larger and larger rescaled annuli.  Two difficulties remain: the field
has only an upper bound, and the total rescaled electron mass can diverge.
Exact neutrality makes its spherical mean nonnegative.  The upper bound
then controls the full local $L^1$ norm, so harmonic interior estimates
give compactness even for these signed fields.

The intermediate subsolutions transfer a strictly positive lower bound
to each limit.  Dilation compactness and contact at approximate extrema
then identify the field as the homogeneous Sommerfeld profile, without
assuming radiality.  The compactness and rigidity criteria are stated
separately in \cref{prof:compactness,prof:rigidity}.

Finally, a conditional profile must be related to the expected mass
defining \eqref{intro:radius}.  A localized $L^2$ count estimate removes
the exceptional observations; shrinking packet supports remove the
smoothing; and a separate dyadic estimate controls the distant tail.
The resulting mass asymptotic, proved along every admissible sequence,
gives both iterated limits by a finite-threshold argument in $Z$.

The proof follows these steps in order.  Section~\ref{sec:scales}
selects observations on expanding annuli, Section~\ref{prof:section}
identifies the limiting field, and Section~\ref{sec:tails} passes to
expected raw counts and proves the radius limits.

\subsection{Notation}

Write $K(x)=|x|^{-1}$, $V=ZK$, and $f(t)=4\pi k(t_+)^{3/2}$.
Constants $C,c>0$ may change from line to line; their permitted
dependencies will be specified.
Expectations include the spin sum and any stated auxiliary randomness.
For the raw spatial configuration of $\Psi_Z$, antisymmetry gives
\begin{equation}\label{intro:count-density}
 \E\#\{i:x_i\in A\}=\int_A\rho_{\Psi_Z}(x)\dd x.
\end{equation}
All finite-system radial boundary spheres have probability zero.
